\chapter{AWS Selection Rationale and Ecosystem}
\label{chap:aws_overview}

\section{Why Amazon Web Services (AWS)?}
Amazon Web Services (AWS) was selected as the cloud infrastructure provider for this system due to its comprehensive service ecosystem, industry-leading reliability, fine-grained identity security, and cost-effective free-tier options for academic prototypes. 

Key technical factors supporting the selection of AWS include:
\begin{itemize}
    \item \textbf{Deep Integration Ecosystem}: AWS services (S3, Lambda, API Gateway, RDS, Cognito, SNS, CloudWatch, IAM, VPC) feature native, API-level integration. IAM roles and security policies allow services to authorize inter-service requests securely without embedding hardcoded access keys in application source code.
    \item \textbf{Serverless Maturity}: AWS Lambda and Amazon API Gateway offer mature execution environments with sub-second scaling, automatic concurrency management, and robust integration with VPC private subnets.
    \item \textbf{Managed PostgreSQL Relational Engine}: Amazon RDS for PostgreSQL eliminates administrative database management by handling automated backups, point-in-time recovery, multi-AZ failover, and minor version patching.
    \item \textbf{Enterprise-Grade Security Standards}: AWS complies with global standards including ISO 27001, SOC 1/2/3, and HIPAA, providing built-in AES-256 encryption at rest and TLS 1.3 encryption in transit.
\end{itemize}

\section{AWS Regional Architecture}
The system deployment is anchored in the \texttt{us-east-1} (N. Virginia) AWS Region. The regional infrastructure distributes workloads across multiple independent Availability Zones (AZs). Each AZ comprises one or more discrete physical data centers with redundant power, networking, and connectivity.

By deploying Amazon RDS in a Multi-AZ configuration across private subnets in \texttt{us-east-1a} and \texttt{us-east-1b}, and deploying AWS Lambda across both subnets, the system maintains continuous operation even if an entire physical data center experiences a local power or network outage.
